\documentclass[10pt,french]{letter}
\usepackage[T1]{fontenc}
\usepackage[utf8]{inputenc}
\usepackage[french]{babel}
\usepackage{a4}
\usepackage{epsfig}
\usepackage{amssymb}
\usepackage{amsmath}

\usepackage{fancyhdr}
\usepackage{hyperref}
\usepackage[left=3cm,right=3cm,top=3cm,bottom=3cm]{geometry}
\usepackage[breakable, skins]{tcolorbox}
\usepackage{enumitem}

% ── Reviewer section header ───────────────────────────────────────────────────
\newcounter{remarkcount}

\newcommand{\reviewer}[1]{%
  \setcounter{remarkcount}{1}%
  \medskip
  \noindent\textbf{\large Reviewer #1}
  \smallskip\par
  \noindent\rule{\linewidth}{0.6pt}
  \medskip
}

% ── Remark box (gray) — numérotation automatique, reset par \reviewer ─────────
\newtcolorbox{remark}{
  breakable, enhanced,
  before={\stepcounter{remarkcount}},
  colback  = gray!12,
  colframe = gray!55,
  fontupper= \small\itshape,
  title    = {\small\bfseries\upshape Remark~\theremarkcount},
  attach boxed title to top left={yshift=-2mm, xshift=4mm},
  boxed title style={colback=gray!55, colframe=gray!55},
  left=3pt, right=3pt, top=5pt, bottom=3pt,
}

% ── Response box (teal left bar) ─────────────────────────────────────────────
\newtcolorbox{response}{
  breakable, enhanced,
  colback  = teal!5,
  colframe = teal!60,
  leftrule = 3pt, rightrule=0.4pt, toprule=0.4pt, bottomrule=0.4pt,
  fontupper= \small,
  title    = {\small\bfseries Response},
  attach boxed title to top left={yshift=-2mm, xshift=4mm},
  boxed title style={colback=teal!60, colframe=teal!60},
  left=6pt, right=3pt, top=5pt, bottom=3pt,
}

\pagestyle{fancy}

\cfoot{\tiny Laboratoire de Mathématiques, UMR CNRS 5127 --- UFR SceM, campus scientifique --- 73376 Le-Bourget-du-Lac Cedex }

\renewcommand{\headrulewidth}{0.4pt}
\renewcommand{\footrulewidth}{0.4pt}


\def\NAME{Romain Negro}
\address{\NAME \\
Laboratoire de Mathématiques, UMR CNRS 5127\\
UFR SceM, Campus scientifique \\
73376 Le-Bourget-du-Lac Cedex\\
}
\signature{\NAME}
\date{Le-Bourget-du-Lac, le \today}

\pagestyle{fancy}

\begin{document}

  \mbox{~}\hfill\mbox{Le-Bourget-du-Lac, April 25, 2026}
  \medskip

  Dear editors,

  Please find with the letter the revised paper (one with the changes highlighted, one without).

  We would like to thank the reviewers for their thorough reading of our manuscript and their
  constructive remarks, which helped us improve the quality of the paper significantly.
  We have carefully addressed all the comments and concerns raised, and provide below a
  detailed account of the changes made in response to each remark.

  \reviewer{1}

  \begin{remark}
    Usually the visibility problem defines two points as visible if and only if the straight-line
    segment joining them lies completely in the free (empty) space. Your definition of visibility
    (page 1, right column, bottom) defines two points as visible if such line segment lies close
    enough to set X. This seems contradictory with the usual definition (and with the intuition)
    unless X is the set of empty voxels (but the reader expects that it is the set of occupied voxels)
    or you address a different problem wrt the classic one. In any case, this needs to be clarified.
    If the addressed problem is not clear, the overall paper cannot be understood properly.
  \end{remark}
  \begin{response}
    You can picture the classic definition of visibility applied to our work as follows: instead
    of considering the surface of the object as the surface of a volume, we consider the surface
    as a thin layer of voxels (like a shell) and we ask for the visibility to be defined as the
    existence of a straight-line segment joining two points that lies completely in this thin layer.
    We added a figure to better illustrate this point (Figure 3 in the revised version) and linked
    it to the definition of visibility in the paper. We hope this will clarify the definition of
    visibility we are using in our work.
  \end{response}

  \begin{remark}
    The second paragraph on page 2, right column, cannot be understood in this place, as the reader
    does not know how the proposed method works. It can be maybe better moved to the conclusions.
  \end{remark}
  \begin{response}
    We moved this paragraph to the conclusion section, as suggested by the reviewer.
  \end{response}

  \begin{remark}
    The paper needs a table of used symbols. Many symbols are used once, without saying that in the rest of
    the paper such symbol will always denote such entity, and then used several lines or pages later, where
    the reader has forgotten what they denote. Moreover, the use of symbols is not always consistent.
  \end{remark}
  \begin{remark}
    The digital set is sometimes denoted with X and other times with Z (for example on page 3, left column,
    both symbols are used). I suggest not to use Z because it is too similar to the set of integer numbers.
    Symbol K seems to be used with the same meaning as X, Z in the right column of page 3.
  \end{remark}
  \begin{remark}
    Symbol h is used on page 2 (right column, top), on page 4 (right column, bottom) and only on page 6 it
    is defined as the gridstep (not in main text, but within the caption of fig.6).
  \end{remark}
  \begin{response}
    Considering how many symbols are used in the paper, we think that a table of symbols would be too heavy
    and would not be really helpful for the reader. However, we added a notation section at the start of the
    visibility section, where we introduce the main symbols used in the paper and explain what they denote.
    We also made sure to be consistent with the use of symbols throughout the paper.
  \end{response}

  \begin{remark}
    The figures, connecting the points of the digital set with bold line segments (like figs. 1 and 2),
    are misleading. They lead to think that X is a 1D complex consisting of points and edges, whereas
    it is a set of points.
  \end{remark}
  \begin{remark}
    The second paragraph of the left column of page 4 seems contraddictory. It says that the input for
    visibility is a subcomplex, while previously it was a set of points in digital space.
  \end{remark}
  \begin{response}
    There was indeed an inconsistency in the figures: the bold line segments connecting lattice
    points made the digital set look like a 1D complex of points and edges, rather than a plain set
    of points. We removed those segments from Figures 1 and 2, and now display only the lattice
    points together with the 1-cells of the cubical complex to make the distinction clear.

    Regarding the apparent contradiction raised in Remark~\the\numexpr\value{remarkcount}-1\relax:
    there is no inconsistency between "a set of points in digital space" and "a subcomplex". The
    key is the \emph{lattice map} (Khalimsky encoding), which assigns to every cell of the cubical
    complex $\mathcal{C}^d$ a unique point with integer coordinates in $\mathbb{Z}^d$. Through this
    bijection, any cell subcomplex (in particular the $k$-star $\mathrm{Star}_k(Z)$ of the digital
    set $Z$) can equivalently be described as a set of points in $\mathbb{Z}^d$.
  \end{response}


  \begin{remark}
    The lattice map (page 4) must be explained better. Figure 3 does not help, unless it is
    explained in more detail in the text. The caption talks about 63 cells, but I cannot
    understand what cells are they in the figure. The process of passing from the cell complex
    to the intervals must be explained.
  \end{remark}
  \begin{response}
    We took inspiration from the original paper's DGMM presentation of the lattice map, and added
    a more detailed explanation of how the lattice map works and how it allows to pass from the cell
    complex to the intervals' representation. We hope this will clarify the concept of lattice map
    and how it is used in our work.
  \end{response}

  \begin{remark}
    Figure 4 (the working example of the algorithm) must be explained in the text, otherwise it is
    useless. What are Vi , Fi , Ti and Ri appearing in the figure? The algorithm is one of the main
    contributions of the paper, so the reader must be enabled to understand it.
  \end{remark}
  \begin{response}
    We added a much more detailed explanation of the working example of the algorithm, including a
    step-by-step description of what is happening in the figure and what the different sets Vi, Fi,
    Ti and Ri represent.
  \end{response}

  \begin{remark}
    Figure 5: the source point (now blue on grey) should be marked in another, more evident, color,
    for example red. Does this figure show 1-visibility?
  \end{remark}
  \begin{response}
    We redesigned Figure 5 to make the source points more visible (now in red and slightly bigger),
    and set the visible points in yellow on newer images. We also extended the visibility radius to
    its maximum value to avoid any confusion about what visibility is being shown in the figure.
    The figure do show the 1-visibility as noted in the visibility introductory text ("Unless otherwise
    specified, we will consider visibility with k = 1").
  \end{response}

  \begin{remark}
    Page 7, top of left column: the visibility normal is the last eigenvector... The normal should be
    a unit vector. Do you mean that it is the eigenvector is taken after normalization, or do you
    guarantee that it 1 has norm = 1? Please clarify.
  \end{remark}
  \begin{response}
    By definition, the eigenvectors of a covariance matrix are unit vectors, so there is no need to
    normalize them, nor to precise that they are unit vectors.
  \end{response}

  \begin{remark}
    Page 7, top of right column: convergence is proven starting from k = 2. Do you mean for any
    k $\geq$ 2 or for any k $\leq$ 2?
  \end{remark}
  \begin{response}
    We meant for any k $\geq$ 2, as stated in the convergence theorem (Section 4) and implied
    by the wording. We clarified this point in the text using the notation k $\geq$ 2 instead of k = 2.
  \end{response}

  \begin{remark}
    Figure 13: I note that the first three rows and columns present vertical artifacts on some faces,
    while in the rightmost column (new method) they don’t appear. I think it is worth to investigate
    why. It is a good feature of the proposed method to avoid such artifacts!
  \end{remark}
  \begin{response}
    Those vertical artifacts are not the only artifacts present in the first three rows and columns
    of Figure 13, but they are indeed more visible than the others. We added a paragraph in the
    section discussing the results of Figure 13 to explain that those artifacts are due to the fact that
    radii of II and PCA normal estimation methods lack the ability to adapt to the local geometry (either
    being too small and capturing inherent digital noise, or being too large and smoothing out salient features),
    while our method is able to adapt to the local geometry and thus avoid such artifacts. We also added two
    versions of radii for these methods in the figure to illustrate this point: one with a small radius (above
    the diagonal) and one with a large radius (below the diagonal).
  \end{response}

  \begin{remark}
    Page 10, top of right column. You should explain what does it mean that the boundary has
    positive reach (what is the reach?). Here X denotes a compact subset of the real space, it
    denoted a digital set of points elsewhere.
  \end{remark}
  \begin{response}
    The \emph{reach} of a set $\partial X$ (introduced by Federer, 1959) is the largest $\rho > 0$
    such that every point within distance $\rho$ of $\partial X$ has a unique nearest point on
    $\partial X$. Intuitively, it measures how "round" the boundary is: a sphere of radius $r$
    has reach $r$, while a sharp corner has reach $0$. The condition $\mathrm{reach}(\partial X) > \rho > 0$
    therefore ensures that the boundary is smooth enough (no sharp features) for the convergence
    proof to hold.
    We also clarified the notation issue concerned with $X$ and $Z$ in the paper,
    and made sure to use $X$ consistently to denote the compact subset of $\mathbb{Z}^d$ and $Z$ to
    denote the digital sets.
  \end{response}

  \begin{remark}
    I did not follow the proof because I don’t have the necessary mathematical knowledge. I hope
    that the other reviewers can cover this part. The term ”sandwiched” (page 11, right column,
    line 5 from bottom) should be probably replaced by a more technical one. Page 13, top: what
    does it mean ”at most” orthogonal? In the caption of Fig. 14, ”1st” should be ”first”.
  \end{remark}
  \begin{response}
    The term "sandwiched" sounds indeed a bit informal, but it is a common term in mathematics
    to describe a situation where an object is "trapped" between two other objects. In our case, the
    estimated normal vector is "sandwiched" between the true normal vector and the tangent space of the
    surface, meaning that it is close to both of them. We fixed the several typos you pointed out.
  \end{response}

  \reviewer{2}


  \begin{remark}
    There are many places where it could have been more precise. For instance:
    "there is a straight segment joining them that stays in the close vicinity of
    the surface" - what is the relevance of this close vicinity and the visibility?
  \end{remark}
  \begin{response}
    
  \end{response}

  \begin{remark}
  \end{remark}
  \begin{response}
  \end{response}
  \begin{remark}
  \end{remark}
  \begin{response}
  \end{response}



  \begin{remark}
  \end{remark}
  \begin{response}
  \end{response}

  We hope this revision and the above answers will satisfy all your
  requests and concerns, and are looking forward to hearing from you soon.

  Best regards
  \bigskip

  \mbox{~}
  \hfill
  \begin{minipage}{0.32\textwidth}
    Romain Negro \\
    PhD in Computer Science\\
    University Savoie Mont Blanc\\
    CNRS, LAMA
  \end{minipage}
  \hfill
  \begin{minipage}{0.32\textwidth}
    Jacques-Olivier Lachaud\\
    Professor in Computer Science\\
    University Savoie Mont Blanc\\
    CNRS, LAMA
  \end{minipage}
  \hfill
  \mbox{~}
\end{document}

